\chapter{Learning Framework, Observation Design, and Reward Formulation}
\label{ch:learning_framework}
\label{ch:perception}

\section{Observation Modalities}

The simulator supports three observation families so that perception can be studied as an ablation rather than a fixed design choice:
\begin{itemize}
    \item \textbf{State-only}: an 8-dimensional vector containing steering angle, hitch angle, tractor cross-track and heading error, trailer cross-track and heading error, and two local curvature terms.
    \item \textbf{State + lidar}: the same 8-dimensional state vector augmented with $N$ normalised range readings, where the beam count $N \in \{8,16,24,32\}$ is varied in the lidar ablation of Chapter~\ref{ch:results}.
    \item \textbf{State + BEV image}: a dictionary observation with the state vector and a $32 \times 32$ grayscale bird's-eye-view image.
\end{itemize}

This design allows the thesis to compare hand-engineered geometric sensing against learned visual representations within the same tractor-trailer task. The three families form a deliberate progression---from a fully engineered state vector, through a mid-level geometric range sensor, to raw bird's-eye-view imagery---so that the marginal value of richer perception can be measured against its added representation-learning and inference cost rather than assumed.

\section{Bird's-Eye View Observation}

The BEV image is generated by first rendering the normal colour Pygame scene and then applying an observation-only transformation before the image is passed to the policy. The base renderer is therefore unchanged for human visualisation, while the policy receives a compact grayscale crop derived from the rendered frame. The current BEV observation is a $32 \times 32$ single-channel image. It is obtained from a local $20\,\mathrm{m}$ square crop of the transformed BEV frame: the forward observation uses the top portion of the local frame to bias the crop toward upcoming path geometry, while the reverse observation uses the bottom portion to bias the crop toward the region behind the vehicle. The image encodes lane markings, obstacle positions, local path geometry, and vehicle geometry, providing spatial context that the vector observation alone cannot capture.

Forward and reverse variants of the BEV environment share the same image size and channel format. They differ only in the crop direction used before resizing and in the interpretation of the vector component. In reverse manoeuvres, the vector component is interpreted with trailer-led tracking errors because the trailer becomes the leading unit during backward motion.

\section{CNN Encoder}

A Convolutional Neural Network (CNN)~\cite{lecun1998gradient} serves as the front-end of the RL agent, processing the BEV image to extract spatial features relevant to lane tracking and obstacle avoidance. The CNN enables the agent to perceive the environment directly from rendered images, without requiring a pre-defined symbolic representation of every object.

The CNN architecture is implemented as a custom feature extractor within the Stable Baselines3 framework~\cite{raffin2021sb3}. The current end-to-end image branch is deliberately small, matching the scale of the $32 \times 32$ observation: a two-layer convolutional encoder with $1 \rightarrow 5 \rightarrow 5$ channels maps the image to an 80-dimensional feature vector, followed by a linear projection to 64 image features. A separate two-layer multilayer perceptron embeds the 8-dimensional vector state to 64 dimensions. The two branches are concatenated before being passed to the TD3 actor and critic networks, producing a 128-dimensional feature vector for the policy and critic heads. Actor and critic feature extractors are not shared, following the TD3 default, so that critic representation updates do not directly overwrite the actor's feature representation.

The CNN can be trained end-to-end with the TD3 agent, allowing the feature extractor to learn which visual features are most relevant for path maintenance and collision avoidance. Initial experiments showed that inserting a randomly initialised CNN branch into an otherwise successful state-based policy can harm early TD3 learning: the agent converged to a short-horizon wall-collision behaviour even though the same vector state was sufficient for rapid learning. To isolate this effect, two diagnostic BEV encoders are included. The \texttt{state\_only} BEV encoder keeps the dictionary observation and multi-input policy path but ignores the image branch, confirming that the BEV environment and multi-input wrapper can still learn from the vector state. The \texttt{scaled\_cnn} encoder introduces the image branch gradually using a scalar schedule:
\begin{equation}
    \phi(s) =
    \left[
    \lambda_I(k)\,\phi_I(I),\,
    \phi_x(\tilde{x})
    \right],
    \label{eq:scaled_cnn_features}
\end{equation}
where $\phi_I$ is the image encoder, $\phi_x$ is the vector-state encoder, and $\lambda_I(k)$ ramps from zero to one over training. A second schedule can reduce the non-steering vector components after the image branch has reached full weight:
\begin{equation}
    \tilde{x} =
    [\delta,\,
    \lambda_x(k)\gamma,\,
    \lambda_x(k)e_y,\,
    \lambda_x(k)e_\psi,\,
    \lambda_x(k)e_{y,t},\,
    \lambda_x(k)e_{\psi,t},\,
    \lambda_x(k)\kappa_1,\,
    \lambda_x(k)\kappa_2].
    \label{eq:state_scale_features}
\end{equation}
The steering angle $\delta$ is intentionally preserved because it is not reliably inferable from the $32 \times 32$ BEV image.

\section{Autoencoder Latent Space}

The diagnostics above motivate decoupling representation learning from control. If a randomly initialised image branch can derail early TD3 learning even when the underlying vector state is easily learnable, then supplying the agent with a usable visual representation from the first training step should remove that failure mode. To improve RL training efficiency, an autoencoder can be used to pretrain an image representation for high-dimensional BEV observations before policy optimisation~\cite{lange2010deep}. The encoder $E(\cdot)$ maps the image observation $s$ to a feature vector $z$, and the decoder $D(\cdot)$ reconstructs the original image:

\begin{equation}
    z = E(s)
    \label{eq:encoder}
\end{equation}
\begin{equation}
    \hat{s} = D(z)
    \label{eq:decoder}
\end{equation}

The autoencoder is pre-trained on images collected from simulation rollouts, allowing the encoder to learn compact representations of the road geometry and obstacle configurations. BEV images are collected from a mixture of simple pure-pursuit rollouts, already-trained compatible policies, and random actions to increase state coverage during pretraining. After reconstruction training, the encoder weights are transferred into the RL feature extractor and can be frozen or fine-tuned while the policy is trained. This produces an ablation between end-to-end image learning, fixed latent representations learned before RL, and pretrained representations that continue adapting during control learning.

\section{UNet Encoder Variant}

In addition to the CNN autoencoder, a UNet-style~\cite{ronneberger2015unet} autoencoder is used for BEV representation learning. The encoder uses repeated double-convolution blocks with max pooling, followed by global average pooling to produce a compact 128-dimensional image descriptor. During pretraining, the full UNet decoder retains skip connections to improve image reconstruction quality; for RL, only the encoder path is kept and concatenated with the 64-dimensional vector embedding.

This second pretrained encoder provides a stronger structural prior than the plain CNN autoencoder. It is intended to test whether preserving multi-scale spatial detail during representation learning improves downstream control performance, particularly in scenes where obstacle shape and lane boundaries must both be interpreted from the BEV image.

\section{Lidar Observation}

The obstacle environments expose a simulated planar lidar observation by appending normalised obstacle distances to the state vector. In forward driving the lidar pose is attached to the tractor, while in reverse driving it is attached to the trailer and rotated by $\pi$ so that the sensing direction follows the leading unit. This design keeps the sensing model consistent with the physical task: the front of the articulated system changes when the direction of travel reverses.

\section{Observation Space Summary}

The full observation interface used across experiments is summarised in Table~\ref{tab:observation_components}.

\begin{table}[H]
\centering
\caption{Observation space components}
\label{tab:observation_components}
\begin{tabular}{p{3cm} p{5cm} p{5cm}}
\toprule
Component & Description & Role \\
\midrule
\texttt{vector} & 8-dimensional state: $[\delta,\gamma,e_y,e_\psi,e_{y,t},e_{\psi,t},\kappa_1,\kappa_2]$ & Provides precise articulation and path-tracking state information \\
\texttt{lidar} & $N$ normalised obstacle ranges, with $N \in \{8,16,24,32\}$ & Encodes local free space around the leading unit without image processing \\
\texttt{image}  & $32 \times 32$ grayscale BEV crop derived from the Pygame render & Encodes lane layout, local path shape, obstacle positions, and vehicle geometry \\
\bottomrule
\end{tabular}
\end{table}

Collision detection is performed using Pygame sprite mask overlap, enabling pixel-precise detection of contact between the vehicle body and obstacles or lane boundaries. Episode termination is triggered on collision, jackknife ($|\gamma| > 90^\circ$), or path completion.

% TODO: architecture diagrams, UNet description, BEV camera anchor configuration, attention visualisation

\section{Twin Delayed Deep Deterministic Policy Gradient (TD3)}
\label{ch:rl_framework}

Twin Delayed Deep Deterministic Policy Gradient (TD3) is an off-policy actor-critic algorithm designed for continuous action spaces. TD3 extends deterministic policy-gradient methods with three key improvements: twin critic networks (to reduce overestimation bias), delayed policy updates (to improve stability), and target policy smoothing (to reduce variance in critic updates). These properties make TD3 a natural fit for articulated-vehicle control: the action space (steering rate and, where enabled, speed) is continuous, a deterministic policy avoids the action-distribution spread that can aggravate the non-minimum-phase reverse dynamics, and the off-policy replay buffer amortises the comparatively high cost of image-based rollouts by reusing every transition many times.

The TD3 actor network $\mu(s \mid \theta^\mu)$ maps states to deterministic actions, while the critic networks $Q_1, Q_2(s, a \mid \theta^{Q_1}, \theta^{Q_2})$ evaluate action quality. The target value is computed using the minimum of the two critic estimates:

\begin{equation}
    y = r_t + \gamma_d \min_{i=1,2} Q_i(s_{t+1},\, \mu(s_{t+1}) + \epsilon)
    \label{eq:td3_target}
\end{equation}

where $\epsilon \sim \text{clip}(\mathcal{N}(0, \sigma), -c, c)$ is clipped Gaussian noise added to the target action for smoothing. Parameters are updated by minimising the temporal difference error in the critics and maximising the expected return through the actor using the deterministic policy gradient theorem~\cite{silver2014dpg}.

TD3 is implemented using the Stable Baselines3 library~\cite{raffin2021sb3}, which provides a reliable implementation integrated with the Gymnasium environment API~\cite{towers2024gymnasium}.

\subsection{Reward Function}
The thesis evaluates several reward variants rather than a single scalar shaping rule. Let
$e_y$ and $e_\psi$ denote the tractor cross-track and heading errors, $e_{y,t}$ and
$e_{\psi,t}$ the trailer cross-track and heading errors, $\gamma$ the hitch articulation
angle, $\dot{x}$ the longitudinal speed, and $\delta$ the commanded steering action. The
guide controller action is denoted $\delta_{\mathrm{pp}}$, with steering normalization
$\Delta_\delta = (\delta - \delta_{\mathrm{pp}})/(2\delta_{\max})$.
The lane-following base environment also defines a proximity penalty
$P_{\mathrm{prox}}$ that is subtracted from non-terminal step rewards. This
penalty is zero when the vehicle is farther than $d_{\mathrm{prox}}=2\,\mathrm{m}$
from blocked occupancy-grid cells and increases quadratically to
$P_{\max}=5$ as clearance approaches zero:
\begin{equation}
P_{\mathrm{prox}} =
\begin{cases}
0, & c \ge d_{\mathrm{prox}},\\
P_{\max}\left(1-\frac{c}{d_{\mathrm{prox}}}\right)^2, & c < d_{\mathrm{prox}},
\end{cases}
\end{equation}
where $c$ is the minimum clearance between the rendered vehicle footprint and
the lane-boundary/obstacle occupancy grid. In the no-obstacle lane-following
environment this discourages driving near lane boundaries; obstacle environments
inherit the same term for both lane boundaries and inserted obstacle cells.

\paragraph{Forward lane-following rewards.}
For forward driving, the terminal reward is
\begin{equation}
R_{\mathrm{term}}^{\mathrm{fwd}} =
\begin{cases}
100, & \text{success}, \\
-100, & \text{failure},
\end{cases}
\end{equation}
and the common shaping terms are
\begin{align}
r_p^{\mathrm{fwd}} &= 0.5\,\dot{x}, \\
P_{\mathrm{tractor}} &= e_y^2 + e_\psi^2, \\
P_{\mathrm{trailer}} &= (0.5\,e_{y,t})^2 + (0.5\,e_{\psi,t})^2.
\end{align}
The four forward reward variants are then
\begin{align}
R_{\mathrm{dense}}^{\mathrm{fwd}} &= r_p^{\mathrm{fwd}} - P_{\mathrm{tractor}} - P_{\mathrm{trailer}} - P_{\mathrm{prox}},
\\
R_{\mathrm{tractor}}^{\mathrm{fwd}} &= r_p^{\mathrm{fwd}} - P_{\mathrm{tractor}} - P_{\mathrm{prox}},
\\
R_{\mathrm{mult}}^{\mathrm{fwd}} &= 4\exp\!\left(-|e_y|-0.5|e_\psi|-0.75|e_{y,t}|-0.5|e_{\psi,t}|\right)\exp\!\left(-1.5|\gamma|\right) \nonumber\\
&\quad + 0.5\,\mathrm{clip}(\dot{x}, 0, 1) - P_{\mathrm{prox}},
\\
R_{\mathrm{guide}}^{\mathrm{fwd}}(k) &= \alpha_k
\left(r_p^{\mathrm{fwd}} - 5\Delta_\delta^2\right)
+ (1-\alpha_k)\left(r_p^{\mathrm{fwd}} - P_{\mathrm{tractor}} - P_{\mathrm{trailer}}\right)
- P_{\mathrm{prox}},
\end{align}
with curriculum weight
\begin{equation}
\alpha_k = \max\!\left(0, 1-\frac{k}{T_{\mathrm{guide}}}\right),
\end{equation}
where $k$ is the current training step in the guided transition schedule and
$T_{\mathrm{guide}}$ is the total number of transition timesteps.

\paragraph{Reverse lane-following rewards.}
For reverse driving, the terminal reward is
\begin{equation}
R_{\mathrm{term}}^{\mathrm{rev}} =
\begin{cases}
200, & \text{success}, \\
-500, & \text{failure},
\end{cases}
\end{equation}
with shared shaping terms
\begin{align}
r_p^{\mathrm{rev}} &= 3\,\mathrm{clip}(-\dot{x}, 0, 1), \\
P_{\mathrm{path}} &= e_y^2 + 0.5\,e_\psi^2 + 0.5\,e_{y,t}^2 + 0.25\,e_{\psi,t}^2,
\\
P_{\mathrm{jack}} &= 10\max(0, |\gamma| - 0.4)^2.
\end{align}
The reverse variants are
\begin{align}
R_{\mathrm{dense}}^{\mathrm{rev}} &= r_p^{\mathrm{rev}} - P_{\mathrm{path}} - P_{\mathrm{jack}} - P_{\mathrm{prox}},
\\
R_{\mathrm{no\mbox{-}hitch}}^{\mathrm{rev}} &= r_p^{\mathrm{rev}} - P_{\mathrm{path}} - P_{\mathrm{prox}},
\\
R_{\mathrm{mult}}^{\mathrm{rev}} &= 5\,\mathrm{clip}(-\dot{x}, 0, 1)\exp\!\left(-|e_y|-0.5|e_\psi|-0.75|e_{y,t}|-0.5|e_{\psi,t}|\right) \nonumber\\
&\quad \times \exp\!\left(-2|\gamma|\right) - P_{\mathrm{prox}},
\\
R_{\mathrm{guide}}^{\mathrm{rev}}(k) &= \alpha_k
\left(r_p^{\mathrm{rev}} - 5\Delta_\delta^2 - P_{\mathrm{jack}}\right)
+ (1-\alpha_k)\left(r_p^{\mathrm{rev}} - P_{\mathrm{path}} - P_{\mathrm{jack}}\right)
- P_{\mathrm{prox}}.
\end{align}

\paragraph{Obstacle-aware extensions.}
Obstacle environments inherit the underlying forward or reverse reward, including
the proximity penalty above, and add a difficulty-dependent speed bonus during
non-terminal episodes:
\begin{equation}
R_{\mathrm{obs,step}} = R_{\mathrm{base}} + d\left(1-\frac{|\dot{x}|}{\dot{x}_{\max}}\right)c_{\mathrm{slow}},
\end{equation}
where $d \in [0,1]$ is obstacle difficulty and $c_{\mathrm{slow}} = 0.1$. If an obstacle
episode terminates through failure before success, the reward is replaced by
\begin{align}
s(d) &= \frac{e^{k_d d}-1}{e^{k_d}-1}, \qquad k_d = 3, \\
R_{\mathrm{fail}}^{\mathrm{obs}} &= p\,R_{\max}\,s(d),
\end{align}
where $p \in [0,1]$ is progress along the route and $R_{\max}=50$ for forward obstacle
tasks and $R_{\max}=100$ for reverse obstacle tasks. If the vehicle instead performs a
deliberate stop in front of an obstacle, the terminal reward becomes
\begin{equation}
R_{\mathrm{stop}}^{\mathrm{obs}} = 1.5\,p\,R_{\max}\,s(d).
\end{equation}


\subsection{Multiplicative vs.\ Additive Reward Structures}

Table~\ref{tab:reward_ablation_forward} and Table~\ref{tab:reward_ablation_reverse} compare these reward variants directly. The dense and tractor/no-hitch forms are additive penalties, while the multiplicative form forces the agent to reduce tracking and articulation errors jointly instead of collecting partial credit from only one term. The guided reward is a curriculum-shaped interpolation: it initially rewards agreement with a pure-pursuit controller~\cite{coulter1992pure} and then decays smoothly into the dense objective.

\section{Hyperparameter Optimisation}

The TD3 agent requires careful hyperparameter tuning. The actor-critic learning rates, replay buffer size, target network update rate, and noise scale are selected over simulation rollouts. For fixed-speed path-tracking experiments the action space is one-dimensional and contains only steering-rate control; when speed remains controllable, exploration noise is kept larger on the speed dimension than on steering. The training pipeline supports matched agents across state-only, lidar, and BEV observation variants. BEV policies use custom multi-input feature extractors, while the state-only and lidar policies use standard multilayer perceptron policies.

To reduce wall-clock cost for image-based training, the training script supports vectorised environment execution with subprocess workers. TD3 optimisation is performed with step-based updates so that the replay buffer and gradient schedule remain responsive for both single-environment and vectorised runs. Vectorised execution is still important for the CNN-based BEV policy because it improves throughput by overlapping environment interaction across subprocess workers, but it no longer changes the underlying update schedule.

\subsection{Representation Pretraining}

Two of the BEV variants use an explicit pretraining stage before TD3 optimisation. The autoencoder-based model transfers a compact pretrained CNN encoder, while the UNet-based model transfers a pretrained UNet encoder after reconstruction training. Both pretraining procedures operate on BEV images collected from policy rollouts and simple geometric controllers. This separates representation learning from control learning and allows direct comparison against the end-to-end CNN baseline. A second set of BEV training diagnostics uses feature scaling: the image branch can be introduced gradually after a state-based policy has begun to solve the task, and the non-steering state inputs can later be reduced to test whether the image representation can replace the engineered tracking state.

\subsection{Classical Controller Tuning}

A reproducible tuning pipeline is used for the classical baselines in later comparison chapters. Forward and reverse pure-pursuit, PID, and MPC controllers are optimised using differential evolution~\cite{storn1997differential} over repeated simulation episodes. The scalar objective penalises trailer cross-track error, jackknife events, collisions, and non-completion. After optimisation, each controller is validated on held-out seeds and a two-dimensional sensitivity sweep is generated over the most influential parameters. This ensures that comparisons against TD3 are made against tuned baselines rather than hand-picked nominal gains.

\section{Failure Replay Curriculum}
\label{sec:failure_replay}

A persistent challenge in training RL agents on the reverse lane-following task is that the episode failure rate is high early in training and the lane geometry is re-randomised at every reset. In the standard random-reset regime, each failure is discarded and replaced by a fresh, randomly drawn layout. As a result, the agent rarely encounters the same hard configuration twice, reducing the pressure to master any particular geometry and potentially slowing convergence on the inherently open-loop-unstable reverse task.

To address this, a \emph{failure replay curriculum} is implemented as a transparent environment wrapper. At each episode boundary the wrapper inspects the episode outcome: if the episode terminated due to a failure---jackknife, collision, or out-of-bounds departure before reaching the goal---the next reset reuses the same random seed that generated the current path and spawn position, exactly replaying the same scenario. If the episode was completed successfully or timed out, a fresh seed is drawn and a new random layout is used. A configurable consecutive-retry ceiling (default five retries) prevents the agent from being trapped indefinitely on a pathologically difficult layout.

This strategy is a form of automatic curriculum learning~\cite{bengio2009curriculum}: rather than constructing a manually graded task sequence, the curriculum emerges from the agent's own failure signal. Its emphasis on repeatedly revisiting the transitions the agent handles worst is shared, in spirit, with prioritized experience replay~\cite{schaul2016per}, though here the emphasis acts at the level of episode sampling rather than replay-buffer weighting. It is implemented orthogonally to the reward formulation and observation representation ablations. The wrapper modifies only the episode-level sampling policy---not the per-step reward signal or the information available to the agent---and is therefore controlled by a separate training flag. Models trained with and without failure replay are saved to distinct directories to ensure clean comparison.

The curriculum is most meaningful for scenarios where the per-episode failure rate is high and layout difficulty is non-uniform, namely reverse lane-following and reverse obstacle avoidance. For forward lane-following the failure rate is low after minimal training and the benefit is expected to be negligible.

% Add final TD3 hyperparameter table here if a dedicated sweep is completed.
